\documentclass[11pt]{article}
%

\usepackage{amsmath, amssymb, amsthm, enumerate}
\usepackage{graphicx, caption, subcaption}
\usepackage{xcolor}
\usepackage{cancel} %
\usepackage[pagebackref]{hyperref}

\usepackage{xspace}


\usepackage{tikz}
\usetikzlibrary{arrows.meta}

%
\definecolor{diagramred}{RGB}{178,32,45}
\definecolor{diagramgreen}{RGB}{0,108,72}
\definecolor{diagramblue}{RGB}{166,206,227}




%
%

%
\usepackage[margin=1.5in, marginparwidth=1.3in]{geometry}

%
\usepackage{titlesec}
\setcounter{secnumdepth}{3}

%


\theoremstyle{plain}
\newtheorem{thm}{Theorem}[section]
\newtheorem{prop}[thm]{Proposition}
\newtheorem{lemma}[thm]{Lemma}
\newtheorem{sublemma}[thm]{Sublemma}
\newtheorem{claim}[thm]{Claim}
\newtheorem{cor}[thm]{Corollary}
\newtheorem{fact}[thm]{Fact}
\newtheorem{conj}[thm]{Conjecture}


\theoremstyle{definition}
\newtheorem{defn}[thm]{Definition}
\newtheorem{assumption}[thm]{Assumption}
\newtheorem{operation}[thm]{Operation}
\newtheorem{construction}[thm]{Construction}
\newtheorem{question}[thm]{Question}
\newtheorem{problem}[thm]{Problem}


%
\newtheorem{rmkx}[thm]{Remark}
\newenvironment{rmk}
{\pushQED{\qed}\renewcommand{\qedsymbol}{$\triangle$}\rmkx}
{\popQED\endrmkx}

\newtheorem{examplex}[thm]{Example}
\newenvironment{example}
{\pushQED{\qed}\renewcommand{\qedsymbol}{$\triangle$}\examplex}
{\popQED\endexamplex}

%
%
%
\newenvironment{proofclaim}[1][Proof]{\begin{proof}[#1]\renewcommand*{\qedsymbol}{\(\square\)}}{\end{proof}}


%
\renewcommand\qedsymbol{$\blacksquare$} 

\newcommand{\barH}{\overline{\mathcal{H}}}
\newcommand{\barPH}{\overline{\mathbb{P}\mathcal{H}}}
\newcommand{\barMg}{\overline{\mathcal{M}}_g}
\newcommand{\barMgn}{\overline{\mathcal{M}}_{g,n}}
\newcommand{\barM}{\overline{\mathcal{M}}}
\newcommand{\C}{\mathbb{C}}
\newcommand{\Z}{\mathbb{Z}}
\newcommand{\R}{\mathbb{R}}
\newcommand{\Q}{\mathbb{Q}}
\newcommand{\E}{\mathbb{E}}
\newcommand\supp{\operatorname{supp}}

\newcommand{\M}{\mathcal{M}}
\newcommand{\calM}{\mathcal{M}}
\newcommand{\calN}{\mathcal{N}}
\newcommand{\K}{\mathcal{K}}
\newcommand{\T}{\mathcal{T}}
\newcommand{\calT}{\mathcal{T}}
\newcommand{\calQ}{\mathcal{Q}}

\newcommand\ninf{n \to \infty}

\renewcommand{\H}{\mathcal{H}}
\renewcommand{\P}{\mathbb{P}}
\renewcommand{\Re}{\operatorname{Re}}
\renewcommand{\Im}{\operatorname{Im}}
\newcommand{\lcm}{\operatorname{lcm}}
\newcommand{\tr}{\operatorname{Trace}}
\newcommand{\area}{\operatorname{area}}
\newcommand{\vol}{\operatorname{vol}}
%

\newcommand\RR{\mathbb{R}}
\newcommand\QQ{\mathbb{Q}}
\newcommand\CC{\mathbb{C}}

\newcommand{\Mgn}{\overline{\mathcal{M}}_{g,n}}
\newcommand\St[1][(\mu)]{\mathcal{H}#1}
%

\newcommand{\TM}{Teichm\"uller\xspace}
\newcommand{\osc}{\operatorname{osc}}



%

\newcounter{bencomments}
\newcommand{\ben}[1]{\textbf{\color{purple}(B\arabic{bencomments})}\marginpar{\textbf{\color{purple}(B\arabic{bencomments})Ben:}
#1}\addtocounter{bencomments}{1}}

 \newcommand{\fred}[1]{\textbf{\color{red} Fred: #1}}

 \newcommand{\john}[1]{\textbf{\color{red} John: #1}}


 
%
%
%
%
%


%
\begin{document}
\begin{problem} 
  Let $\M_g$ be the moduli space of genus $g$ Riemann surfaces.  We consider the Deligne-Mumford compactification $\overline\M_g = \M_g \sqcup \partial \M_g$. 
 The boundary $\partial \M_g$ admits a partition into strata  $\Delta_{G}$ parametrizing stable curves with dual graph $G$.   

 Let $V_0 \subset \M_g$ be a (complex) smooth subvariety of (complex) dimension $2$ that is totally geodesic with respect to the Teichm\"uller metric.  Let $V_t$ be a holomorphic family of subvarieties of $\M_g$ deforming $V_0$, for $t$ in the unit disc $\Delta \subset \C$ 
 \footnote{i.e. a complex submanifold $\mathcal V \subset \M_g \times \Delta$, with $V_t = \psi(\phi^{-1}(t))$, where $\phi: \mathcal V \to \Delta, \psi: \mathcal V \to \M_g$ are the natural projections; we will further assume that this gives a smoothly trivial family of submanifolds. }.
 We impose several technical assumptions on this family, namely (all closures below, e.g. $\overline V_t$, are taken in $\overline \M_g$):
 \begin{enumerate}[(i)]
 \item \label{item:smooth} for every boundary stratum $\Delta_{G}$, we have that $\overline V_t \cap \Delta_{G}$ is a (smoothly trivial) family, and %
   
 \item \label{item:compatible} the families so produced are compatible with $V_0$ in the sense that if the boundary of $V_0$ in a smaller stratum arises from its boundary in a larger stratum, then the same holds for $V_t$.  That is, if $\Delta_G, \Delta_{G'}$ are strata with $\Delta_{G'}$ on the boundary of $\Delta_G$, and $U$ is an open set in $\overline \M_g$ with 
   \begin{align*}
        \overline{\overline V_0 \cap \Delta_{G}} \cap \Delta_{G'} \cap U  = \overline V_0 \cap \Delta_{G'} \cap U , 
   \end{align*}
   %
   then, for all $t$,
      \begin{align*}
        \overline{\overline V_t \cap \Delta_{G}} \cap \Delta_{G'} \cap U= \overline V_t \cap \Delta_{G'} \cap U.
      \end{align*}

 \end{enumerate}
%
 Note that we do \emph{not} assume that $V_t$ is totally geodesic for $t\ne 0$.  

Now fix $\Gamma$ a dual graph; we will assume for ease of notation that $\Delta_\Gamma$ is a product of lower complexity moduli spaces $\M_{g_1,n_1} \times \cdots \times \M_{g_k,n_k}$ \footnote{for some boundary strata, we would need to quotient this by a finite group; we exclude those here}.  For each $i$, let $\pi_i: \Delta_\Gamma \to \M_{g_i,n_i}$ be the natural projection map.

Denote by $\partial V_t$ the closure of $V_t$ in $\overline{\M} _g$ intersected with $\Delta_\Gamma$.  Suppose that $\dim \partial V_t=1$ and that each $\partial V_t$ is constant on some factor $\M_{g_i,n_i}$ of $\Delta_\Gamma$, i.e. for some $i$, $\pi_i(\partial V_t)=\{Y_t\}$ for all $t$.   Prove that the family $\partial V_t$ is non-varying on that factor, i.e. $\pi_i(\partial V_t) = \{Y_0\}$ for all $t$.   

   %
  
  %
\end{problem}
\end{document}
